%% ============================================================
%% shared/services/compute/cicd_deployment.tex
%% CI/CD, Deployment Strategies & SSM — SAP (some SAA)
%% ============================================================

\servicesection{CI/CD Pipeline Services}{\saptag}{CodeCommit, CodeBuild, CodeDeploy, CodePipeline}

\begin{keyfacts}
\begin{itemize}
  \item \textbf{CodeCommit}: managed Git repository; private, encrypted, integrated with IAM.
        (Note: AWS announced end-of-life for new customers in 2024 --- GitHub/GitLab often preferred.)
  \item \textbf{CodeBuild}: fully managed build service; compiles source, runs tests, produces artifacts.
        Pay per build minute. No servers to manage.
  \item \textbf{CodeDeploy}: automates application deployments to EC2, on-premises servers, Lambda,
        and ECS. Supports rolling, blue/green, and canary deployment strategies.
  \item \textbf{CodePipeline}: orchestrates the full CI/CD workflow; connects source $\rightarrow$
        build $\rightarrow$ test $\rightarrow$ deploy stages. Integrates with GitHub, CodeCommit,
        CodeBuild, CodeDeploy, CloudFormation, Elastic Beanstalk, and ECS.
  \item \textbf{CodeArtifact}: managed artifact repository for npm, Maven, PyPI, NuGet packages.
        Proxies public package registries with a security boundary.
\end{itemize}
\end{keyfacts}

\begin{examtip}
The standard SAP CI/CD pipeline is: \textbf{Source} (CodeCommit / GitHub)
$\rightarrow$ \textbf{Build} (CodeBuild)
$\rightarrow$ \textbf{Deploy} (CodeDeploy / CloudFormation / ECS)
orchestrated by \textbf{CodePipeline}.
Each stage can trigger an approval action (manual gate) before proceeding.
\end{examtip}

\subsection{Deployment Strategies}

\begin{comparebox}[title={Deployment strategy comparison}]
\begin{tabular}{@{}p{3.4cm}p{4.2cm}p{1.5cm}p{2.4cm}@{}}
\toprule
\textbf{Strategy} & \textbf{How it works} & \textbf{Downtime} & \textbf{Rollback} \\
\midrule
\textbf{All-at-once (In-place)} & Deploy to all instances simultaneously &
  Yes (brief) & Re-deploy old version \\[4pt]
\textbf{Rolling}        & Deploy to a subset at a time; old and new coexist &
  Minimal & Re-deploy old \\[4pt]
\textbf{Rolling + extra batch} & Launch extra capacity first; then roll; maintains capacity &
  None & Re-deploy old \\[4pt]
\textbf{Immutable}      & Launch brand-new instances; replace old entirely &
  None & Terminate new \\[4pt]
\textbf{Blue/Green}     & New environment in parallel; switch traffic (DNS or ALB) &
  None & Flip DNS back \\[4pt]
\textbf{Canary}         & Shift small \% of traffic to new version; monitor; increase &
  None & Route 100\% back \\
\bottomrule
\end{tabular}
\end{comparebox}

\begin{examtip}
\textbf{Blue/Green} is the preferred SAP answer for ``zero-downtime'' + ``easy rollback''.
Control traffic shifting with Route~53 weighted routing or ALB target group weights.
CodeDeploy natively supports blue/green for EC2, Lambda (traffic shift), and ECS.
\textbf{Canary} (e.g.\ 10\% $\rightarrow$ 100\%) is best for gradual validation with automated rollback
on CloudWatch alarm.
\end{examtip}

\begin{sapdepth}
\textbf{SAP: Advanced deployment patterns:}
\begin{itemize}
  \item \textbf{Lambda Canary/Linear}: CodeDeploy shifts Lambda traffic in configurable steps
        (e.g.\ \texttt{LambdaCanary10Percent5Minutes}: 10\% for 5 min, then 100\%). Automatic rollback
        on CloudWatch alarm.
  \item \textbf{ECS Blue/Green}: CodeDeploy registers a new task definition; shifts ALB traffic
        between blue (old) and green (new) target groups. Original task set is retained for rollback window.
  \item \textbf{CloudFormation Change Sets}: preview infra changes (like a diff) before applying.
        Use in CodePipeline for IaC deployments with human approval gates.
  \item \textbf{Service Catalog}: package CloudFormation templates as approved Products;
        developers self-service without direct CF access. Enforces governance at scale.
  \item \textbf{AWS Proton}: managed deployment service for container and serverless applications;
        platform teams define environment/service templates; dev teams self-deploy.
\end{itemize}
\end{sapdepth}

%% ────────────────────────────────────────────────────────────────

\servicesection{AWS Systems Manager (SSM)}{\bothtag}{Unified operational management for EC2 and hybrid infrastructure}

\begin{keyfacts}
\begin{itemize}
  \item \textbf{Agentless or agent-based}: SSM Agent pre-installed on Amazon Linux 2 / Windows AMIs.
        On-premises servers use the agent with a hybrid activation.
  \item No need to open SSH (port 22) or RDP (port 3389) for management operations.
  \item SSM communicates over HTTPS to the SSM endpoint (no inbound ports required).
\end{itemize}
\end{keyfacts}

\subsection{Key SSM Capabilities}

\begin{comparebox}[title={SSM features}]
\begin{tabular}{@{}p{3.8cm}p{8cm}@{}}
\toprule
\textbf{Feature} & \textbf{Purpose} \\
\midrule
\textbf{Session Manager}        & Interactive shell access to EC2 and on-premises without SSH/RDP.
                                  Session logs go to S3/CloudWatch. Replaces bastion hosts. \\[4pt]
\textbf{Parameter Store}        & Hierarchical key-value store for config and secrets.
                                  Standard tier free; Advanced tier supports larger values and policies (TTL). \\[4pt]
\textbf{Run Command}            & Execute commands at scale across a fleet of instances without SSH.
                                  Output goes to S3/CloudWatch. \\[4pt]
\textbf{Patch Manager}          & Define patch baselines (which patches are approved); schedule
                                  maintenance windows; scan and remediate. \\[4pt]
\textbf{State Manager}          & Define and maintain desired configuration state of instances
                                  (e.g.\ antivirus installed, port 22 closed). \\[4pt]
\textbf{Inventory}              & Collect metadata: installed software, OS config, network config.
                                  Query across fleet; feed into Config. \\[4pt]
\textbf{Automation}             & Runbooks (YAML/JSON) for automated remediation workflows
                                  (e.g.\ restart service, create snapshot, rotate key). \\[4pt]
\textbf{OpsCenter}              & Central location for operational issues (OpsItems); integrates
                                  with CloudWatch alarms, Config, GuardDuty findings. \\
\bottomrule
\end{tabular}
\end{comparebox}

\begin{comparebox}[title={Parameter Store vs Secrets Manager}]
\begin{tabular}{@{}p{3.5cm}p{4cm}p{4cm}@{}}
\toprule
\textbf{Feature} & \textbf{Parameter Store} & \textbf{Secrets Manager} \\
\midrule
Cost            & Free (Standard) / \$0.05/10k API (Advanced) & \$0.40/secret/month + API \\
Auto-rotation   & No native rotation & Yes (built-in for RDS/Aurora) \\
Max size        & 4 KB (Standard) / 8 KB (Advanced) & 64 KB \\
Cross-account   & Via IAM / resource policy & Native cross-account \\
Versioning      & Yes & Yes \\
Use for         & Config values, non-sensitive params, simple secrets & DB passwords, API keys, auto-rotated credentials \\
\bottomrule
\end{tabular}
\end{comparebox}

\begin{examtip}
\textbf{``Automatically rotate database credentials''} $\Rightarrow$ \textbf{Secrets Manager}.
\textbf{``Store application configuration cheaply''} $\Rightarrow$ \textbf{Parameter Store} (Standard tier is free).
\textbf{``Replace bastion host''} $\Rightarrow$ \textbf{SSM Session Manager} (no open ports, full audit trail).
\end{examtip}
